\scanpage{206}
Considérons le cas où $t'\in S'{}^+$ est un élément régulier de $S'$ ; alors $\lambda$ tel que $u(t)=\lambda t'$ est uniquement déterminé par $u$ ; donc si on se borne aux
triples $(S,S^+\ \text{à p.d.},t)$ tels que $t$ soit régulier, et aux hom $u:(S,S^+\deleted{,t})\longrightarrow(S',S'{}^+\deleted{,t'})$ tels que $u(t)\in S.t'$, alors $A_\ast(S,S^+,t)$ devient un foncteur sur cette catégorie, à valeurs ds la cat. des
\deleted{\unsure{alg}} anneaux gradués \inserted{(alternés)} différentiels \deleted{\unsure{act|et}} filtrés.

La condition envisagée sur $(S,S^+_{\mathrm{pd}},t)$ est vérifiée en particulier si on fait $S=k\{T\}$, $S^+=(T)_{\mathrm{pd}}$, $t=T$, lorsque $k$ est sans torsion.

Revenons alors à $\displaystyle A_\ast(k)\overset{\mathrm{dfn}}{=}A_\ast(\underset{\substack{\Vert\\k\{T\}}}{S_k},S_k^+,T)$, et à la détermination des endom de $A_\ast(k)$, comme \deleted{anneau} \inserted{$k$-algèbre} différentielle $\frac12$ simpliciale (où on ne considère \deleted{\unsure{donc}} la structure de $S$-algèbre, ni la graduation, ni les p.d.\deleted{\unsure{?}}, ni de filtration a priori\ldots)
\deleted{\unsure{Bien entendu|Bien entendu que|lorsqu'on} \unsure{que si}} Soit $\unsure{G|\mathscr G}_k$ le \deleted{groupe} \inserted{monoïde} des couples $(u,\lambda)$, où $u$ est un endom. de $S$ comme $k$-algèbre à p. div., et $\lambda\in S^{\deleted{\unsure{\times|+}}}$ tel que $u(T)=\lambda\deleted{\{T\}}T$.
Comme $u\deleted{\unsure{(T)}}$ est connu par la connaissance de $u(T)$,
$G_k$ est donc $\simeq$ le \deleted{groupe} \inserted{monoïde}
\deleted{\unsure{des} $S^\times$} \inserted{d'ens. sous-jacent $S$, avec la multiplication $(\lambda'\ast\lambda)\{T\}\deleted{\unsure{T}}=\lambda'\{T\}\deleted{\unsure{T}}\{T\lambda'\{\unsure{\lambda|T}\}\}$}. On désigne par $(u_\lambda)_\ast$ l'endomorphisme de $A_\ast(k)$ défini par $\lambda\in S$ (d'où $u_\lambda u_{\lambda'}=u_{\lambda\lambda'}$)

Il est clair que $\lambda\longmapsto u_{\lambda\ast}$ est un hom de monoïdes \uline{injectif} (en fait $\lambda\longmapsto(u_\lambda)_{\unsure{1|k|\ast}}$ l'est déjà, car $(u_\lambda)_\ast(X_0)=\lambda X_0$ et $A_1(k)\simeq k\{T\}\{X_0\}$. Ceci dit,
Si $\unsure{H|\mathscr H}_k$ est le monoïde des $G_k$-endomorphismes de la $k$-algèbre différentielle $\frac12$ simpliciale $A_\ast(k)$, \deleted{est-il vrai que} $\Omega_k$ le monoïde de \uline{tous} les endomorphismes de la $k$-alg. diff. $\frac12$ simpliciale, on a
% archive sheet number: 205
